\scanpage{189}
Le calcul ``bébête'', \unsure{en développant|du développ.t} le déterminant et tenant compte des symétries (pour regrouper les termes qui se correspondent par \deleted{symétrie} op. de $\unsure{\mathfrak S_3|\mathscr S_3}$) et $\Sigma^2$,
\[
 \Sigma^2-\Delta^2=4\Pi(b,b',b'',t,t',t'',c,c',c'')
\]
avec
\[
\boxed{\begin{aligned}
 \Pi={}&tt't''+(ct^2+c't^{\prime 2}+c''t^{\prime\prime 2})+(b^2c'c''+b^{\prime 2}c''c+b^{\prime\prime 2}cc')\\
 &-(b'b''ct+b''bc't'+bb'c''t'')
\end{aligned}}
\]
\[
 =tt't''+ct(t-b'b'')+c't'(t'-b''b)+c''t''(t''-bb')+(b^2c'c''+b^{\prime 2}c''c+b^{\prime\prime 2}cc')
\]

On aura donc
\[
 \left\{\begin{aligned}
 \tau+\tau'&=\Sigma&&\bigl(=bt+b't'+b''t''-bb'b''\bigr)\\
 \tau\tau'&=\Pi&&\bigl(=tt't''+(ct^{\unsure{2|\prime}}+c't^{\prime 2}+c''t^{\prime\prime 2})^{\unsure{,|\prime}}\\
 &&&\quad+(b^2c'c''+b^{\prime 2}c''c+b^{\prime\prime 2}cc')\\
 &&&\quad-(b'b''ct+b''bc't'+bb'c''t'')\bigr)
 \end{aligned}\right.
\]
\[
 \tau-\tau'=\varepsilon_0\Delta
\]
% archive sheet number: 188
